% Generated draft: requires mathematical and source review.
\documentclass[11pt]{article}
\usepackage[margin=25mm]{geometry}
\usepackage{fontspec}
\setmainfont{Latin Modern Roman}
\usepackage{amsmath,amssymb,mathtools,unicode-math}
\setmathfont{Latin Modern Math}
\usepackage{enumitem,needspace,xurl,hyperref}
\setlength{\parindent}{0pt}
\setlength{\parskip}{5pt}
\setlength{\emergencystretch}{3em}
\newcommand{\N}{\mathbb{N}}
\newcommand{\Z}{\mathbb{Z}}
\newcommand{\Q}{\mathbb{Q}}
\newcommand{\R}{\mathbb{R}}
\newcommand{\C}{\mathbb{C}}
\newcommand{\F}{\mathbb{F}}
\newcommand{\A}{\mathbb{A}}
\newcommand{\PP}{\mathbb{P}}
\newcommand{\E}{\mathbb{E}}
\newcommand{\eps}{\varepsilon}
\newcommand{\dd}{\,\mathrm{d}}
\begin{document}
\section*{UnsolvedMath ID: 30004967}
\textbf{Source problem code:} catalogue entry
\subsection{Definitions and mathematical statement}
The supplied record does not determine a precise mathematical problem. It refers to an unidentified Theorem~3, omits its hypotheses and conclusion, and omits the subsequent results and corollary. Consequently, no quantified mathematical statement can be reconstructed from this record.

The excerpt describes the Colin de Verdi\`ere parameter \(\mu(G)\) as the largest corank of certain matrices associated with a graph \(G\). For an \(n\times n\) matrix \(A\), its corank is \(n-\operatorname{rank}(A)\). However, the admissible matrices and their defining conditions are absent, so the supplied description does not constitute a definition of \(\mu(G)\). The permitted class of graphs is also unspecified. The final sentence begins a characterization involving \(\mu(G)\leq 1\) and terminates at ``if and''; its missing equivalence must not be completed by inference.

The excerpt asserts that proper minor-closed graph families have bounded expansion and that Theorem~3 yields a result answering a question of Eppstein~[1], a corollary disproving a conjecture of Huynh and Wood~[4], and another result answering a question of Huynh and Wood~[4]. None of these questions, the conjecture, or the resulting conclusions is supplied. Their relationship to the subsequent discussion of \(\mu(G)\) cannot be established from the excerpt.

The registry labels the record's status as \texttt{open}, while its statement status is \texttt{unrecoverable}. Its supplied verification note reports that the official report contains a theorem and corollary about the Colin de Verdi\`ere parameter rather than an open problem, and that no separate clean statement was promoted. These are supplied registry assessments, not independently verified mathematical conclusions. This statement-only draft records the unresolved formulation; it does not formulate a recoverable conjecture.
\subsection{Short English statement}
The source fragment discusses consequences of an unidentified theorem and then introduces an incompletely defined graph parameter. The theorem, questions, conjecture, corollary, and final characterization are missing. A precise mathematical target cannot be recovered, and the supplied \texttt{open} label does not establish that an identifiable problem remains open.
\subsection{Sources}
\begin{itemize}
\item Local registry: \texttt{data/unsolvedmath/problems.json}, permanent UnsolvedMath ID \(30004967\), exact source code \texttt{catalogue entry}, title ``Unspecified Colin de Verdiere Graph Parameter Problem,'' in the Oberwolfach Reports Open Problems collection. The supplied record retains the original excerpt for provenance and supplies no clean statement.
\item Supplied source citation: \emph{Graph Theory (2023)}, Oberwolfach Reports, DOI \texttt{10.4171/owr/2022/1}, pp.~5--78. Supplied source URL: \url{https://doi.org/10.4171/owr/2022/1}. The background describes Oberwolfach Reports as reports of mathematical workshops and their open-problem discussions. The citation's 2023 label differs from the 2022 report designation and proposed year; this discrepancy is unresolved.
\item Supplied literature source: ``Oberwolfach Report 1/2022,'' \url{https://doi.org/10.4171/owr/2022/1}. The literature assessment dated exactly \texttt{2026-08-22} labels the record open but states that it is only a meta-level description of an unnamed theorem resolving or refuting earlier questions, without an identifiable mathematical claim whose final status can be checked. This is supplied dated curation derived from the OWR record, not independently verified; time-sensitive status and cited claims require independent verification.
\item Attribution within the excerpt: Eppstein~[1] and Huynh and Wood~[4]. The corresponding bibliography entries, questions, and conjecture are absent. No proposer is identified, and no authoritative citation specifying a recoverable target problem is supplied.
\end{itemize}
\end{document}
